\documentclass[11pt,a4paper]{article}
\usepackage[margin=1in]{geometry}
\usepackage{amsmath,amssymb,amsthm}
\usepackage{booktabs}
\usepackage{array}
\usepackage[T1]{fontenc}
\usepackage{lmodern}
\usepackage{microtype}
\usepackage{textcomp}
\usepackage{xcolor}
\definecolor{seclink}{HTML}{1F4E8C}
\usepackage[colorlinks=true,linkcolor=seclink,citecolor=seclink,urlcolor=blue]{hyperref}
\usepackage{parskip}
\usepackage[most]{tcolorbox}

\definecolor{ink}{HTML}{1A1A1A}
\definecolor{grey}{HTML}{5F5B54}
\definecolor{yellowbg}{HTML}{FBF3D5}
\definecolor{pinkbg}{HTML}{FADDE1}

\tcbset{plainbg/.style={boxrule=0pt, sharp corners, colframe=white,
  left=4mm, right=4mm, top=3mm, bottom=3mm, before skip=3mm, after skip=3mm}}
\newtcolorbox{sectionbox}{plainbg, colback=yellowbg}
\newtcolorbox{alarmbox}{plainbg, colback=pinkbg}

\newtheorem{claim}{Claim}
\newtheorem*{obs}{Observation}

\title{\vspace{-2em}A Sheaf Theory of Federation}
\author{\normalsize Local agreement, global obstruction, and where the deadlock actually lives}
\date{\normalsize Claude, 7 September 2026}

\begin{document}
\maketitle
\vspace{-1em}

\begin{sectionbox}
The historical note argued that a federation is a diagram of morphisms whose graph
has cycles, and that the four combinators build trees. That is a statement about
what cannot be composed. It says nothing about how badly a given federation fails,
and it offers no invariant.

Sheaf theory supplies both. The claim of this note is that a federation is the
assertion that a certain presheaf is a sheaf, that the obstruction to that
assertion is a cohomology class, and that the class is carried by the fundamental
group of the federation. The cycles were the right thing to worry about. What was
missing was that they act.
\end{sectionbox}

\section{The site, and the presheaf}

Let $X$ be the polity and let $\{U_i\}$ be an open cover, one $U_i$ per member.
An overlap $U_i \cap U_j$ is a matter falling under both jurisdictions: a river
between two provinces, a citizen who moves, trade crossing a boundary.

Define a presheaf $\mathcal{F}$ by taking $\mathcal{F}(U)$ to be the decisions,
rules or legal states valid on $U$, with restriction maps
$\mathcal{F}(U) \to \mathcal{F}(V)$ for $V \subseteq U$ sending a decision to what
it says on the smaller jurisdiction. Restriction along inclusion is the only
structure required, and every federation supplies it, because a rule made for the
whole says something about each part.

\begin{sectionbox}
\textbf{The thesis in one sentence.} A federation is the claim that
$\mathcal{F}$ is a sheaf.
\end{sectionbox}

\section{The two axioms, read politically}

Sheafhood is two conditions and both are political programmes.

\textbf{Separatedness.} If two global decisions agree on every $U_i$ then they are
equal. A federation asserting this is asserting that the centre holds nothing
invisible from below. Its failure is a centre with content no member can see, which
is Gosplan, and which is the subject of \texttt{terror.txt}.

\textbf{Gluing.} If decisions $s_i \in \mathcal{F}(U_i)$ agree on every overlap,
$s_i|_{U_{ij}} = s_j|_{U_{ij}}$, then some global $s$ restricts to all of them.
This is the federation working. Everyone agrees locally, and there exists a
single thing they have agreed on.

\textbf{Failure of gluing is the deadlock, stated without any reference to time.}
Local consistency everywhere, and no global object. Nobody has behaved badly and
there is nothing to decide, because the thing to be decided does not exist.

\begin{obs}
Sheafification is the fourth political option and the notes did not have it. Force
$\mathcal{F}$ into a sheaf and you obtain a global decision, but the sections
change. You get an answer, and it is not the one any member locally held. That is
imposed unification, and it is distinct from all four repairs because it does not
resolve the disagreement. It replaces the disagreeing parties.
\end{obs}

\section{Descent, and the cocycle condition}

Sections are the easy case, since gluing is an axiom for them. The interesting
object is the structure itself. Each member holds not merely a decision but a local
system: a legal order, a currency, a register of persons.

Write $g_{ij}$ for the transition data on $U_{ij}$, the isomorphism identifying
member $j$'s structure with member $i$'s where both apply. A treaty is exactly a
$g_{ij}$. Mutual recognition is exactly a $g_{ij}$.

These glue into a single global structure precisely when, on every triple overlap,
\[
  g_{ij} \circ g_{jk} \;=\; g_{ik} \qquad \text{on } U_{ijk}.
\]

\begin{sectionbox}
\textbf{The cocycle condition is a condition on cycles.} It says that composing
the agreements around a loop returns you to where you started. Bilateral treaties
are never the problem. The problem is that three bilateral agreements can be
individually impeccable and inconsistent as a triangle.
\end{sectionbox}

That is the same structural fact the container argument reached from the other
side. There the cycles obstructed a trace. Here they carry a class.

\section{The obstruction, and where it lives}

Transition data satisfying the cocycle condition, modulo the changes that amount to
relabelling, is exactly $H^1(X, \mathcal{G})$ for $\mathcal{G}$ the sheaf of
automorphisms.

\begin{sectionbox}
$H^0$ is what the federation can jointly decide.

$H^1$ is the set of ways it can be locally consistent and globally twisted.
\end{sectionbox}

For a constant sheaf of abelian groups $G$ on a reasonable connected space,
\[
  H^1(X; G) \;\cong\; \operatorname{Hom}(\pi_1(X), G),
\]
and in the non-abelian case the locally constant classes are
$\operatorname{Hom}(\pi_1(X), G)$ up to conjugacy.

\begin{claim}
The obstruction to federation is a representation of the fundamental group of the
federation.
\end{claim}

This is the sharpest form of the thing the earlier notes were circling. Cycles were
correctly identified as the difficulty. What was missing is that a cycle does not
merely fail to terminate. It \emph{acts}, and the obstruction is what it does.

\section{Monodromy is the right name}

Carry a decision around a loop of members, applying each agreement in turn, and
compare what comes back with what set out. Trivial monodromy means the loop does
nothing and everything glues. Non-trivial monodromy means the federation is
consistent between every pair of members and has no global object.

\begin{sectionbox}
This supersedes the commutativity claim in the historical note, which was correct
and imprecise. \textbf{Commuting operations are the special case of trivial
monodromy.} Monodromy is better because it is not a yes or a no. It is an element
of a group, so it says which twist, by how much, and around which loop.
\end{sectionbox}

\section{The four repairs, in this language}

\begin{center}
\begin{tabular}{@{}l@{\hspace{4mm}}>{\raggedright\arraybackslash}p{0.63\textwidth}@{}}
\toprule
\textbf{Orient} & Pass to the universal cover. On a simply connected base the
locally constant classes vanish, because there is no $\pi_1$ to represent. A
federation without cycles has no obstruction, and this is a theorem rather than a
hope. Bass and Serre give the group-theoretic version: a graph of groups acts on a
tree, and the tree is the universal cover. \textbf{So orienting and unrolling are
one construction.} The centre is available at the price of infinitely many copies
of the federation, which is a precise statement of what a centre costs. \\[2mm]
\textbf{Unroll} & Compute \v{C}ech cohomology at a fixed refinement rather than
passing to the limit. An answer for this cover, not for the space. \\[2mm]
\textbf{Guard} & Guarded recursion already lives in a presheaf topos, the topos of
trees $\mathsf{Set}^{\omega^{\mathrm{op}}}$, where the later modality is an
endofunctor and guarded fixpoints exist. The delay repair was sheaf-theoretic
before we arrived. \\[2mm]
\textbf{Commute} & Trivial monodromy, as above. The cocycle condition holds
automatically and $H^1$ contributes nothing. \\
\bottomrule
\end{tabular}
\end{center}

\section{The toy example, and the historical one}

\textbf{The M\"obius band.} Cover it with two strips. Each member sees a perfectly
ordinary piece of paper with a well defined front and back. The two agree on both
overlaps. There is no global assignment of front and back. Every local view is
correct, every pairwise agreement holds, and the global object does not exist.

That is a federation failing, drawn, with no history attached.

\begin{alarmbox}
\textbf{And now the loaded one, stated cleanly.} Assign to each jurisdiction the
legal status of a person. Each state's assignment is internally consistent. Each
pair of states has a rule for what happens at the boundary. Carry a person around a
loop and the status returned differs from the status that set out.

\emph{Dred Scott} is a monodromy computation. Legal status was a section of a
bundle which was not trivial, the obstruction was the first class of the cover, and
no global assignment existed. The war was not a failure of goodwill and not a
failure of procedure. There was no global section, and the federation had asserted
there was one.
\end{alarmbox}

\section{What is solid and what is modelling}

\begin{alarmbox}
\textbf{Solid.} The sheaf axioms, the cocycle condition, $H^1$ classifying
descent data, the isomorphism with $\operatorname{Hom}(\pi_1, G)$ for locally
constant coefficients, the vanishing on simply connected bases, and the
Bass--Serre theorem. None of that is mine and none of it is in question.

\textbf{Modelling.} Everything upstream of it. Whether decisions on a jurisdiction
genuinely form a presheaf with the right restriction maps, whether jurisdictions
form a site at all rather than merely a set with intersections, and whether the
automorphism sheaf is the right coefficient object. That is where the whole
argument is load bearing, and it is contestable at every step.

\textbf{The question I would put to you.} The container note found that cycles
obstruct a trace on $\triangleleft$. This note finds that cycles carry a class in
$H^1$. If those are the same obstruction, then $H^1$ measures the failure of the
trace to exist, and there ought to be a statement relating them. If they are not
the same, one of the two models is wrong about federation and it would be useful to
know which.
\end{alarmbox}

\end{document}
